% Generated draft: requires mathematical and source review.
\documentclass[11pt]{article}
\usepackage[margin=25mm]{geometry}
\usepackage{fontspec}
\setmainfont{Latin Modern Roman}
\usepackage{amsmath,amssymb,mathtools,unicode-math}
\setmathfont{Latin Modern Math}
\usepackage{enumitem,needspace,xurl,hyperref}
\setlength{\parindent}{0pt}
\setlength{\parskip}{5pt}
\setlength{\emergencystretch}{3em}
\newcommand{\N}{\mathbb{N}}
\newcommand{\Z}{\mathbb{Z}}
\newcommand{\Q}{\mathbb{Q}}
\newcommand{\R}{\mathbb{R}}
\newcommand{\C}{\mathbb{C}}
\newcommand{\F}{\mathbb{F}}
\newcommand{\A}{\mathbb{A}}
\newcommand{\PP}{\mathbb{P}}
\newcommand{\E}{\mathbb{E}}
\newcommand{\eps}{\varepsilon}
\newcommand{\dd}{\,\mathrm{d}}
\begin{document}
\section*{UnsolvedMath ID: 30003859}
\textbf{Source problem code:} catalogue entry
\subsection{Definitions and mathematical statement}
Let \(L\) denote a line configuration that is rigid in the sense intended by the source, and let \(S(n,L)\) denote its Hirzebruch--Kummer covering surface with exponent \(n\). The supplied record does not specify the ambient space or ground field of \(L\), the notion of rigidity for configurations or surfaces, or the construction and smoothness conventions for \(S(n,L)\). These are unresolved qualifications, so the following is a quantified statement of the intended question, conditional on those definitions being fixed.

Using positive integer exponents, the question is whether
\[
\forall L\text{ rigid},\quad
\exists N_L\in\Z_{>0}\quad
\forall n\in\Z_{>0},\qquad
n\geq N_L\ \Longrightarrow\ S(n,L)\text{ is rigid}.
\]
Here the threshold may depend on \(L\). This formulation presupposes that the intended covering surface is defined for every sufficiently large positive integer exponent. Any restrictions on admissible configurations or exponents must be supplied from the cited sources.

The supplied dated curation describes the complete-quadrangle case, denoted \(\mathrm{CQ}\), as solved and reports that \(S(n,\mathrm{CQ})\) is rigid exactly for \(n\geq4\); it describes the assertion for arbitrary rigid line configurations as open. These are supplied literature assessments, not independently verified mathematical claims. Their construction, exponent-domain, and rigidity conventions remain to be checked.
\subsection{Short English statement}
Does every rigid line configuration have rigid Hirzebruch--Kummer covering surfaces for all sufficiently large integer exponents, with the threshold allowed to depend on the configuration? The intended notions of rigidity and the covering construction require clarification. Supplied dated curation labels the problem partially solved and reports a solution for the complete quadrangle.
\subsection{Sources}
\begin{itemize}
\item Local registry: supplied complete source record for permanent UnsolvedMath ID \(30003859\), source problem code \texttt{catalogue entry}, titled \emph{Rigidity of Hirzebruch--Kummer Coverings}, in the Oberwolfach Reports Open Problems collection. The record gives proposed year 2018 and names no individual proposer.
\item Supplied report citation: \emph{Subgroups of Cremona Groups (2019)}, Oberwolfach Reports, pp.~1685--1743, DOI 10.4171/owr/2018/28. \url{https://doi.org/10.4171/owr/2018/28}. The record describes Oberwolfach Reports as documenting mathematical workshops and their open-problem discussions.
\item Bauer--Catanese, \emph{On the rigidity of line arrangements}. \url{https://arxiv.org/abs/1803.02984}.
\item Bauer--Catanese, \emph{Line arrangements and rigidity}. \url{https://arxiv.org/abs/1609.08128}.
\item \emph{Oberwolfach Report 28/2018}. \url{https://ems.press/content/serial-article-files/46750}.
\item Literature assessment dated exactly \texttt{2026-08-21}: the supplied curation labels the problem partially solved, reports rigidity of \(S(n,\mathrm{CQ})\) exactly for \(n\geq4\), and describes the general conjecture as open. This is supplied dated curation derived from the OWR record, not an independently verified assessment; the record explicitly requests independent verification of time-sensitive status and cited claims.
\end{itemize}
\end{document}
